% !TeX root = ../Book.Stress_regulation.tex
\documentclass[../Book.Stress_regulation.tex]{subfiles}
\graphicspath{{\subfix{../images/}}}
\begin{document}

\section{An Example}

Somebody always causes me to be afraid (for instance in the place of work: threats of being fired).
This weakens the lungs.
But a weakness of the lungs also means: a decline of the mental capacities, depression, over sensitivity, emotionally imbalanced, to the point of hysteria.
So, due to the weakness of the lungs the quality of my work suffers. That leads to the threats of my boss getting more pronounced.
I become oversensitive and emotionally unhinged, maybe even depressive or hysteric. So my letter of loosing my job might follow soon!

In this situation of stress it is important to strengthen the lungs (also at work). This would mean in this concrete case:

\begin{itemize}
	\item No more smoking in inside rooms\footnote{I know it's a bit dated.}
	\item Orient on the weakest link in the chain and go in it's pace.
	\item An eventual AC should be switched off (everlasting "wind").
	\item Increasing the turn--over shouldn't be allowed to be an end in itself.
\end{itemize}

Analogous examples could be made with fear of cancer, fear of loneliness, fear of dangers. In any case, organs get weakened and these weaknesses cause in turn other problems.



\epigraph{Everything is linked together, everything is inseparable.}{\textit{Dalai Lama}}
\setlength\epigraphwidth{.4\textwidth} 
\end{document}